\documentclass[11pt]{article}
\usepackage{amsmath,amssymb}
\usepackage{hyperref}
\usepackage{geometry}
\usepackage{algorithm}
\usepackage{algorithmic}
% Use XeLaTeX to compile this file: xelatex online_cn.tex
\usepackage{xeCJK}
\setCJKmainfont{STSong}  % macOS built-in font
\setCJKsansfont{STHeiti}
\setCJKmonofont{STFangsong}
\geometry{margin=1in}

\begin{document}

\section{混合推理的在线算法}
\label{sec:online}

本文档详细介绍混合推理系统的在线路由和调度算法。系统利用两种不同的订阅服务以及一个备用 API：

\begin{enumerate}
  \item \textbf{服务 1 ($S_Q$，如 Chutes)}：受\textbf{每日配额} ($Q$) 约束。
  \item \textbf{服务 2 ($S_C$，如 Featherless)}：受\textbf{瞬时并发限制} ($C$) 约束。
  \item \textbf{备用 API ($S_A$)}：按 token 付费的服务，容量无限但成本较高。
\end{enumerate}

\paragraph{统一路由策略。}
我们采用\textbf{统一路由器}，并行评估将请求路由到 $S_Q$、$S_C$ 或 $S_A$ 的"净价值"，选择最大化预期成本节省减去机会成本的选项。

%==============================================================================
\section{第一部分：服务 1 - 配额限制订阅 ($S_Q$)}
\label{sec:part1}

该组件管理每日配额约束。它作为一个\textbf{自治代理}，跟踪配额使用情况，并向统一路由器提供\textbf{影子价格（阈值）}。

\subsection{问题设定 ($S_Q$)}

\begin{itemize}
  \item \textbf{输入}：多天内到达的请求流 $r_1, r_2, \dots, r_T$。
  \item \textbf{每日配额}：订阅每天可服务的固定请求数 $Q$。配额在每天开始时重置（如 UTC 00:00）。
  \item \textbf{决策}：对于每个请求 $r_t$，决定 $x_t \in \{ \text{订阅}, \text{API} \}$。
  \item \textbf{价值}：将 $r_t$ 路由到订阅的价值是避免的 API 成本：
  \[ v_t = \text{Cost}_{\text{API}}(r_t) \]
  \item \textbf{目标}：在每日配额约束下最大化总价值（节省）：
  \[ \max \sum_{t: x_t = \text{订阅}} v_t \quad \text{s.t.} \quad \sum_{t \in \text{第 } d \text{ 天}} \mathbb{I}(x_t = \text{订阅}) \le Q \]
\end{itemize}

\subsection{基线：贪心策略}

贪心策略是最直观的方法。只要配额可用就使用订阅，不考虑请求的潜在价值。

\paragraph{算法逻辑。}
\begin{enumerate}
  \item \textbf{每日重置}：每天开始时，重置已用配额 $k = 0$。
  \item \textbf{路由}：如果 $k < Q$，路由到\textbf{订阅}并增加 $k$。否则，路由到\textbf{API}。
\end{enumerate}

\paragraph{分析。}
\begin{itemize}
  \item \textbf{优点}：简单；如果需求 $\ge Q$，确保 100\% 利用率。
  \item \textbf{缺点}：\textbf{忽略机会成本}。可能用低价值请求填满配额，迫使高价值请求后来走 API。
\end{itemize}

\subsection{进阶：基于阈值的原始-对偶策略}

为了克服贪心的局限性，我们采用受\textbf{在线背包问题}的\textbf{原始-对偶}方法启发的\textbf{基于阈值}的方法。

\paragraph{核心思想：影子价格。}
我们为配额分配一个动态的\textbf{影子价格}或\textbf{阈值}。该阈值代表请求必须提供的\emph{最小价值}，才能证明消耗一单位稀缺配额是合理的。

\paragraph{阈值函数。}
设：
\begin{itemize}
  \item $Q$：每日总配额
  \item $k$：今天已使用的配额槽数
  \item $z = k/Q$：已使用配额的比例（$z \in [0, 1]$）
  \item $L$：请求价值的估计下界（如第 5 百分位）
  \item $U$：请求价值的估计上界（如第 95 百分位）
\end{itemize}

阈值函数为：
\[ \psi(z) = L \cdot \left( \frac{U}{L} \right)^z \]

\paragraph{在统一路由中的作用。}
对于请求 $r_t$，该组件计算使用 $S_Q$ 的\textbf{机会成本}：
\[ \text{Cost}_{S_Q}(r_t) = \begin{cases} \psi(z) & \text{如果 } k < Q \\ \infty & \text{如果 } k \ge Q \end{cases} \]

\paragraph{不确定性下的成本估计。}
在在线设置中，精确的 API 成本 $v_t$ 是未知的，因为输出 token 数 $n_t^{\text{out}}$ 在生成完成前是未知的。我们使用估计值 $\hat{v}_t$：
\[ \hat{v}_t = \frac{n_t^{\text{in}}}{1000} \cdot p^{\text{in}} + \frac{\hat{n}_t^{\text{out}}}{1000} \cdot p^{\text{out}} \]

估计 $\hat{n}_t^{\text{out}}$ 的策略：
\begin{enumerate}
  \item \textbf{历史平均}：特定模型或数据集的平均输出长度。
  \item \textbf{指数移动平均 (EMA)}：$\hat{n}_{t+1}^{\text{out}} = \alpha \cdot n_t^{\text{actual}} + (1-\alpha) \cdot \hat{n}_t^{\text{out}}$
\end{enumerate}

\paragraph{理论说明。}
对于价值在 $[L, U]$ 范围内的在线背包问题，该算法的竞争比至多为 $\ln(U/L) + 1$，这是渐近最优的。

%==============================================================================
\section{第二部分：服务 2 - 并发限制订阅 ($S_C$)}
\label{sec:part2}

该组件管理并发约束 ($C$)。它作为一个\textbf{自治代理}，管理自己的内部队列和调度，向统一路由器暴露\textbf{拥塞价格}。

\subsection{问题设定 ($S_C$)}

\begin{itemize}
  \item \textbf{输入}：由统一路由器路由到 $S_C$ 的请求。
  \item \textbf{约束}：
    \begin{itemize}
      \item \textbf{总容量}：$C_{\text{total}} = C_{\text{plan}} \times N_{\text{sub}}$
      \item \textbf{当前使用量}：$W_{\text{active}} = \sum_{r \in \text{活跃}} \text{weight}(r)$
      \item \textbf{可用容量}：$C_{\text{avail}} = C_{\text{total}} - W_{\text{active}}$
    \end{itemize}
  \item \textbf{队列容量}：大小为 $K$ 的有限队列缓冲请求。
\end{itemize}

\subsection{处理时间估计 ($p_t$)}

本节估计在并发限制订阅 $S_C$ 上运行请求 $r_t$ 的服务时间。该估计用于 CAPQ 的\textbf{可行性检查}和\textbf{价值密度计算}。

\emph{注意：处理时间 $p_t$ 仅与 $S_C$（并发约束）相关。对于 $S_Q$（仅配额），路由决策仅取决于 API 成本 $v_t$，而非持续时间。}

\paragraph{分解。}
我们将服务时间分解为预填充（首 token 时间）和解码阶段：
\[ p_t \approx \widehat{\text{TTFT}}(m_t) + \frac{\hat{n}^{\text{out}}_t}{\widehat{\text{TPS}}(m_t)} \]

其中：
\begin{itemize}
  \item $m_t$：$r_t$ 请求的模型
  \item $\widehat{\text{TTFT}}(m)$：模型 $m$ 的估计首 token 时间
  \item $\widehat{\text{TPS}}(m)$：模型 $m$ 的估计解码吞吐量（每秒 token 数）
  \item $\hat{n}^{\text{out}}_t$：预测的输出 token 数
\end{itemize}

\paragraph{在线估计器。}
统计数据从最近的 $S_C$ 请求日志计算，按模型分桶：
\begin{itemize}
  \item $\hat{n}^{\text{out}}_t$：来自轻量级预测器或 EMA 基线的预测输出 token。
  \item $\widehat{\text{TTFT}}(m)$：模型 $m$ 观察到的 TTFT 值的稳健统计量（如中位数或 P50）。
  \item $\widehat{\text{TPS}}(m)$：从已完成请求导出：
  \[ \text{TPS} \triangleq \frac{n^{\text{out}}}{\max(\text{延迟} - \text{TTFT}, \epsilon)} \]
\end{itemize}

\paragraph{风险感知估计。}
为避免低估并发占用，我们使用\textbf{置信上界 (UCB)} 估计：
\[ p_t^{\text{UCB}} = \widehat{\text{TTFT}}^{(0.9)}(m_t) + \frac{\hat{n}^{\text{out},(0.9)}_t}{\widehat{\text{TPS}}^{(0.1)}(m_t)} \]

\textbf{解释}：
\begin{itemize}
  \item TTFT 使用\textbf{第 90 百分位}（保守，考虑队列延迟）
  \item 输出长度使用\textbf{第 90 百分位}（考虑长响应）
  \item TPS 使用\textbf{第 10 百分位}（保守，考虑慢解码）
\end{itemize}

\textbf{在 CAPQ 中的使用}：
\begin{itemize}
  \item $p_t^{\text{UCB}}$ 用于计算机会成本：$\text{Cost}_{S_C}(r_t) = \lambda_t \cdot w_t \cdot p_t^{\text{UCB}}$
  \item $p_t^{\text{UCB}}$ 用于价值密度：$\rho_t = v_t / (w_t \cdot p_t^{\text{UCB}})$
\end{itemize}

\subsection{原始-对偶成本感知调度 (CAPQ)}

我们将并发管理表述为带有\textbf{影子价格}的动态资源分配问题。

\paragraph{核心思想：并发影子价格。}
我们使用\textbf{价值密度}为将请求纳入订阅系统定义\textbf{影子价格} $\lambda_t$。

\begin{itemize}
  \item \textbf{情况 1：容量可用}：资源充足。$\lambda_t = 0$。所有请求被接纳。
  \item \textbf{情况 2：系统饱和}：资源稀缺。影子价格为：
  \[ \lambda_t = \min_{r \in \text{队列}} \frac{v(r)}{\text{weight}(r) \cdot p(r)} \]
\end{itemize}

\paragraph{在统一路由中的作用。}
对于请求 $r_t$，该组件计算使用 $S_C$ 的\textbf{机会成本}：
\[ \text{Cost}_{S_C}(r_t) = \lambda_t \cdot \text{weight}(r_t) \cdot p(r_t) \]

%==============================================================================
\section{第三部分：统一路由策略}
\label{sec:part3}

本节定义协调 $S_Q$、$S_C$ 和 $S_A$ 的主逻辑。

\subsection{统一决策逻辑}

对于每个到达的请求 $r_t$：

\paragraph{步骤 1：价值估计。}
\[ v_t = \text{Cost}_{\text{API}}(r_t) \]

\paragraph{步骤 2：获取影子价格（机会成本）。}
\begin{itemize}
  \item $S_Q$：$\theta_Q = \text{Service1.get\_threshold}()$（如配额耗尽则返回 $\infty$）
  \item $S_C$：$\theta_C = \text{Service2.get\_congestion\_price}()$（如模型不兼容则返回 $\infty$）
\end{itemize}

\paragraph{步骤 3：计算净收益。}
\begin{itemize}
  \item \textbf{选项 1 ($S_Q$)}：$G_Q = v_t - \theta_Q$
  \item \textbf{选项 2 ($S_C$)}：$G_C = v_t - \theta_C \cdot \text{weight}(r_t) \cdot p(r_t)$
  \item \textbf{选项 3 ($S_A$)}：$G_A = 0$（基线）
\end{itemize}

\paragraph{步骤 4：决策。}
\[ \text{选择} = \arg\max \{ G_Q, G_C, G_A \} \]

%==============================================================================
\section{第四部分：评估计划}
\label{sec:part4}

\subsection{指标}
\begin{enumerate}
  \item \textbf{总成本节省（\$）}：主要目标（最大化）。
  \item \textbf{节省比率}：$\frac{\text{Savings}_{\text{在线}}}{\text{Savings}_{\text{最优}}} \in [0, 1]$（目标：最大化，越接近 1 越好）。
  \item \textbf{资源利用率}：$S_Q$ 的配额利用率，$S_C$ 的并发利用率。
  \item \textbf{SLO 违规率}：错过截止时间的请求百分比（针对 $S_C$）。
\end{enumerate}

\subsection{实验阶段}

\paragraph{实验 1：每日配额优化 ($S_Q + S_A$)。}
重点：验证原始-对偶阈值相对于贪心的每日配额管理效果。
\begin{itemize}
  \item \textbf{基线}：贪心、离线最优（ILP）、原始-对偶（我们的）
\end{itemize}

\paragraph{实验 2：联合优化 ($S_Q + S_C + S_A$)。}
重点：验证统一路由器和 CAPQ 管理配额和并发的效果。
\begin{itemize}
  \item \textbf{基线}：贪心（优先级 $S_Q \to S_C \to S_A$）、离线最优（完整 ILP）、统一原始-对偶（我们的）
\end{itemize}

%==============================================================================
\section{第五部分：学习增强在线算法}
\label{sec:part5}

本节用\textbf{机器学习预测}扩展原始-对偶框架，以改善不确定性下的路由决策。我们将其形式化为\textbf{学习增强在线算法}，在预测准确时提供一致性保证，在预测失败时提供鲁棒性。

\subsection{动机：信息差距}

在线路由的根本挑战是决策时的\textbf{信息不对称}：

\begin{center}
\begin{tabular}{|l|c|c|}
\hline
\textbf{信息} & \textbf{到达时已知} & \textbf{影响} \\
\hline
输入 token $n^{\text{in}}_t$ & 是 & 较小（成本的 20-30\%）\\
输出 token $n^{\text{out}}_t$ & \textbf{否} & 较大（成本的 70-80\%）\\
处理时间 $p_t$ & \textbf{否} & 对 $S_C$ 关键 \\
\hline
\end{tabular}
\end{center}

API 成本为：
\[ v_t = \frac{n^{\text{in}}_t}{1000} \cdot p^{\text{in}} + \frac{n^{\text{out}}_t}{1000} \cdot p^{\text{out}} \]

由于大多数提供商的 $p^{\text{out}} \approx 3\text{-}4 \times p^{\text{in}}$，准确估计 $n^{\text{out}}_t$ 至关重要。

\subsection{预测框架}

\paragraph{预测目标。}
给定请求特征 $\mathbf{x}_t$（输入长度、模型、提示特征），我们预测：
\begin{enumerate}
  \item \textbf{输出长度估计}：$\hat{L}_t = f(\mathbf{x}_t)$
  \item \textbf{预测不确定性}：$\sigma_t = g(\mathbf{x}_t)$
\end{enumerate}

预测的 API 成本变为：
\[ \hat{v}_t = \frac{n^{\text{in}}_t}{1000} \cdot p^{\text{in}} + \frac{\hat{L}_t}{1000} \cdot p^{\text{out}} \]

\paragraph{通过分位数回归量化不确定性。}
我们使用\textbf{分位数回归}预测输出长度分布的多个分位数，而非点估计：
\begin{itemize}
  \item $\hat{L}^{(10)}_t$：下界预测（第 10 百分位）——用于保守成本估计
  \item $\hat{L}^{(50)}_t$：中位数预测（第 50 百分位）——用于点估计
  \item $\hat{L}^{(90)}_t$：上界预测（第 90 百分位）——用于 $S_C$ 的持续时间估计
\end{itemize}

我们定义相应的成本估计：
\begin{align}
  \hat{v}^{(50)}_t &= \frac{n^{\text{in}}_t \cdot p^{\text{in}} + \hat{L}^{(50)}_t \cdot p^{\text{out}}}{1000} \quad \text{（点估计）} \\
  \hat{v}^{\text{LCB}}_t &= \frac{n^{\text{in}}_t \cdot p^{\text{in}} + \hat{L}^{(10)}_t \cdot p^{\text{out}}}{1000} \quad \text{（保守下界）}
\end{align}

\textbf{解释}：对于配额路由（$S_Q$），我们使用保守的 $\hat{v}^{\text{LCB}}_t$ 以避免在实际价值可能较低的请求上浪费配额。

\paragraph{特征空间。}
我们仅使用请求到达时可用的\textbf{可观测特征}：
\begin{center}
\begin{tabular}{|l|l|}
\hline
\textbf{类别} & \textbf{特征} \\
\hline
请求级别 & 输入 token 数 $n^{\text{in}}_t$、max\_tokens 设置 \\
模型级别 & 模型名称（分类）、模型家族 \\
时间 & 一天中的小时、星期几、会话开始后的时间 \\
历史 & 用户最近输出长度的滚动平均（EMA）\\
\hline
\end{tabular}
\end{center}
\emph{注意：我们不假设能访问提示内容。所有特征都来自请求元数据和历史统计。}

\subsection{风险感知路由算法}

\paragraph{核心公式。}
我们修改原始-对偶决策规则，使用预测成本的\textbf{置信下界 (LCB)}：
\[
x_t =
\begin{cases}
\text{订阅} & \text{如果 } \hat{v}^{\text{LCB}}_t > \psi(z_t) \text{ 且 } k_t < Q \\[5pt]
\text{API} & \text{否则}
\end{cases}
\]

其中：
\begin{itemize}
  \item $\hat{v}^{\text{LCB}}_t$：使用 $\hat{L}^{(10)}_t$（输出长度的第 10 百分位）的保守成本估计
  \item $\psi(z_t) = L \cdot (U/L)^{z_t}$：影子价格函数（阈值）
  \item $z_t = k_t / Q$：当前配额利用率
\end{itemize}

\paragraph{解释。}
通过使用预测输出长度的\textbf{第 10 百分位}，我们获得请求 API 成本的保守（悲观）估计。这确保：
\begin{itemize}
  \item \textbf{高置信度高价值}的请求（即使 P10 也超过阈值）获得配额优先
  \item \textbf{价值不确定}的请求（P10 低尽管 P50 高）被路由到 API
\end{itemize}

这比启发式 $\hat{v} - \lambda \sigma$ 更有原则，因为它直接使用分位数回归输出，无需引入额外的调参参数 $\lambda$。

\paragraph{算法。}

\begin{algorithm}[H]
\caption{学习增强原始-对偶路由}
\begin{algorithmic}[1]
\STATE \textbf{输入：}请求 $r_t$，当前配额使用量 $k_t$，预测器 $\mathcal{M}$
\STATE \textbf{参数：}配额 $Q$，价值边界 $[L, U]$
\STATE
\STATE \COMMENT{步骤 1：分位数预测}
\STATE $\hat{L}^{(10)}_t, \hat{L}^{(50)}_t, \hat{L}^{(90)}_t \gets \mathcal{M}.\text{predict\_quantiles}(r_t.\text{features})$
\STATE
\STATE \COMMENT{步骤 2：计算保守成本估计 (LCB)}
\STATE $\hat{v}^{\text{LCB}}_t \gets (n^{\text{in}}_t \cdot p^{\text{in}} + \hat{L}^{(10)}_t \cdot p^{\text{out}}) / 1000$
\STATE
\STATE \COMMENT{步骤 3：计算影子价格}
\STATE $z_t \gets k_t / Q$
\STATE $\theta_t \gets L \cdot (U/L)^{z_t}$
\STATE
\STATE \COMMENT{步骤 4：风险感知决策}
\IF{$k_t < Q$ \AND $\hat{v}^{\text{LCB}}_t > \theta_t$}
    \STATE $k_{t+1} \gets k_t + 1$
    \RETURN 订阅
\ELSE
    \RETURN API
\ENDIF
\end{algorithmic}
\end{algorithm}

\subsection{理论分析}

\paragraph{学习增强框架。}
我们使用学习增强算法文献中的\textbf{一致性-鲁棒性}框架分析算法。

\textbf{定义（预测误差）：}
\[ \eta_t = |v_t - \hat{v}^{\text{LCB}}_t| \]
其中 $v_t$ 是真实 API 成本，$\hat{v}^{\text{LCB}}_t$ 是保守预测成本。

\textbf{定义（一致性）：}预测准确时（小 $\eta_t$）的节省比率。目标：接近 1。

\textbf{定义（鲁棒性）：}预测任意错误时的最坏情况节省比率。目标：远离 0。

\paragraph{分析。}

\textbf{经验改进。}当预测校准良好时：
\begin{itemize}
  \item LCB 方法比基于 EMA 的估计更有效地过滤低价值请求。
  \item 实验表明，与使用 EMA 的原始原始-对偶相比，该算法实现更接近 1 的节省比率。
\end{itemize}

\textbf{鲁棒性警告。}与已知价值的原始原始-对偶不同，纯分位数路由\textbf{缺乏正式的最坏情况保证}：
\begin{itemize}
  \item 如果预测被对抗性地低估（如 $\hat{L}^{(10)} \approx 0$），算法会将所有请求路由到 API，实现零节省。
  \item 原始-对偶的理论 $O(\ln(U/L))$ 界假设\emph{已知}请求价值；预测误差打破了这一假设。
\end{itemize}

\textbf{实用保障。}为缓解对抗性预测，我们推荐一个后备机制：
\begin{itemize}
  \item 如果 $\hat{v}^{\text{LCB}}_t < \epsilon$（预测可疑地低），回退到基于 EMA 的估计或贪心。
  \item 这种混合方法保留了经验收益，同时避免了灾难性失败模式。
\end{itemize}

\textbf{重要说明。}即使\emph{当前}价值估计完美，在线算法也无法实现节省比率 $\to 1$，因为它们不知道\emph{未来}请求。ML 的好处是减少当前请求价值的估计误差，而不是预测未来到达。

\paragraph{分位数选择。}
分位数的选择（如 P10 vs P20）控制保守程度：

\begin{center}
\begin{tabular}{|c|l|l|}
\hline
\textbf{分位数} & \textbf{行为} & \textbf{权衡} \\
\hline
P10（默认） & 保守 & 可能配额利用不足 \\
P20-P30 & 适中 & 平衡 \\
P50 & 激进 & 可能在不确定请求上浪费配额 \\
\hline
\end{tabular}
\end{center}

\textbf{实用指导}：默认使用 P10；如果配额利用率太低，在验证集上调优。

\subsection{与相关工作的联系}

\paragraph{带预测的滑雪租赁。}
配额路由问题类似于\textbf{滑雪租赁问题}：
\begin{itemize}
  \item 订阅 = 购买（固定成本，配额内无限使用）
  \item API = 租用（按次付费）
\end{itemize}

学习增强滑雪租赁算法（Purohit et al., 2018; Gollapudi \& Panigrahi, 2019）为我们的方法提供了理论基础。

\paragraph{带预测的在线背包。}
我们的问题也可以看作\textbf{在线背包}：
\begin{itemize}
  \item 容量 = 每日配额 $Q$
  \item 物品价值 = API 成本节省 $v_t$
  \item 物品重量 = 1（每个请求消耗一个配额单位）
\end{itemize}

阈值函数 $\psi(z)$ 对应于背包 LP 松弛中的对偶变量。

\subsection{实验设计}

\paragraph{基线。}
\begin{center}
\begin{tabular}{|l|l|}
\hline
\textbf{方法} & \textbf{描述} \\
\hline
贪心 & 如果配额可用就使用订阅（先到先服务）\\
原始-对偶 (EMA) & 基于阈值，使用 EMA 输出长度估计 \\
LA-PD (P50) & 学习增强，使用中位数预测（激进）\\
LA-PD (P10) & 学习增强，使用 P10 预测（保守，默认）\\
离线最优 & 完整信息上界 \\
\hline
\end{tabular}
\end{center}

\paragraph{消融研究。}
\begin{enumerate}
  \item \textbf{预测质量}：向预测注入噪声以测试鲁棒性
  \item \textbf{分位数敏感性}：比较 P10、P20、P30、P50 的成本估计
  \item \textbf{特征重要性}：消融输入特征以识别关键预测器
\end{enumerate}

\paragraph{指标。}

\emph{路由性能：}
\begin{center}
\begin{tabular}{|l|l|}
\hline
\textbf{指标} & \textbf{定义} \\
\hline
节省比率 & $\text{Savings}_{\text{在线}} / \text{Savings}_{\text{最优}} \in [0,1]$ \\
总节省 & $\sum_{t: x_t = \text{订阅}} v_t$（绝对 \$）\\
配额利用率 & 已消耗的每日配额比例 \\
\hline
\end{tabular}
\end{center}

\emph{预测校准：}
\begin{center}
\begin{tabular}{|l|l|}
\hline
\textbf{指标} & \textbf{定义} \\
\hline
分位数覆盖率 & 实际值 $\leq$ 预测分位数的百分比（目标：匹配分位数水平）\\
Pinball 损失 & $\sum_t \rho_\tau (L_t - \hat{L}^{(\tau)}_t)$，其中 $\rho_\tau$ 是分位数损失函数 \\
MAE & $\hat{L}^{(50)}_t$ 的平均绝对误差 \\
\hline
\end{tabular}
\end{center}

\emph{注意：校准指标对于验证风险感知路由至关重要。如果 P10 覆盖率显著 $>10\%$，说明预测过于保守。}

\subsection{总结}

学习增强原始-对偶算法用以下方式扩展了经典的基于阈值的方法：
\begin{enumerate}
  \item \textbf{分位数回归}：预测输出长度的多个分位数（P10、P50、P90）
  \item \textbf{保守估计}：使用下分位数（P10）进行风险感知配额分配
  \item \textbf{实用保障}：当预测可疑地低时回退到 EMA/贪心
\end{enumerate}

该框架桥接了\textbf{在线算法}和\textbf{机器学习}，在实践中实现比原始原始-对偶更高的节省比率。注意，与经典在线算法不同，纯 ML 路由缺乏正式的最坏情况保证；后备机制提供了实用的鲁棒性。

\end{document}
